\documentclass{practical-resume}
%\pagenumbering{arabic} %page numbering

%Custom Variables
\newcommand*{\name}{Rui Sousa} 

% Footer to Direct Readers to Website Resume.
\newcommand{\foot}{Last Update: \today}

%fonts
\setmainfont{TeX Gyre Termes}
\setsansfont{TeX Gyre Heros}

\usepackage{hyperref}
\hypersetup{
    pdftitle={Resume},
    pdfauthor={Rui Sousa},
    colorlinks=true,
    linkcolor=blue,
    citecolor=magenta,
    filecolor=magenta,      
    urlcolor=magenta,
    pdfpagemode=FullScreen
}

\begin{document}

%%%%%%%%%%%%%%%%%%%%%%%%%%%%%%%%%%%%%%%%%%%%%%%%%%%%%%%%%%%%%%%%%%%%%%
% Header
% prints your name
\namehead

% Links
\begin{address}{
	Office 3436, Kellogg Global Hub, \\
	2211 N Campus Dr, Evanston, IL 60208
}
{rui.agm.sousa@u.northwestern.edu}
{https://ruiagmsousa.github.io}
{\today}
\end{address}

%%%%%%%%%%%%%%%%%%%%%%%%%%%%%%%%%%%%%%%%%%%%%%%%%%%%%%%%%%%%%%%%%%%%%%
% Education 
\sectionheading{Education}
\organization{Northwestern University}{Evanston, IL}{2021 -- 2027 (expected)} 
\begin{position}{PhD Candidate, }{Economics Department}
	%\item Master of Arts, 2022  
\end{position} 
\vspace{1em}

\organization{Nova School of Business and Economics}{Lisbon, Portugal}{2016 -- 2019}
	\begin{position}{MSc Economics, }{GPA $19/20$ (A) - Novo Banco Master's Scholarship}
		%\item Exchange Program at University of Copenhagen
	\end{position}
\vspace{0.6em}

%%%%%%%%%%%%%%%%%%%%%%%%%%%%%%%%%%%%%%%%%%%%%%%%%%%%%%%%%%%%%%%%%%%%%%
% Working Papers
\sectionheading{Research Activity}

	\organization{The Aggregate Green Elasticity of Substitution (\href{https://ruiagmsousa.github.io/research/agg-green-elasticity}{Draft})}{}{}
\begin{itemize}[leftmargin=2.5em, itemsep=-0.25em, topsep=0.5em]\vspace{-0.3em}
    \item Develop a semi-structural framework to estimate the short-run aggregate elasticity of substitution between clean and polluting energy sources---a key parameter governing energy price pass-through and the costs of decarbonization
    \item Construct an equilibrium relationship between energy consumption and observable prices, and exploit cross-state variation in energy mixes to implement a shift-share IV strategy, estimating an aggregate elasticity of substitution of $0.5$
    \item Derive and calibrate a mapping between sectoral and aggregate elasticities, finding that end users substitute toward electricity more readily than the electricity sector substitutes toward clean generation ($0.80$ versus $0.52$)
\end{itemize}

    \organization{Cleansing Oil (\href{https://ruiagmsousa.github.io/research/cleansing-oil/}{Draft})}{}{}
\begin{itemize}[leftmargin=2.5em, itemsep=-0.25em, topsep=0.5em]\vspace{-0.3em}
\item Estimate that a 10\% increase in oil prices reduces CO$_2$-equivalent atmospheric concentrations by 0.10 ppm relative to trend
\item Show that petroleum consumption declines across refined products, partially substituted toward low-quality coal and bioenergy
\item Find that oil price shocks reduce solar investment through tighter financial conditions and increased uncertainty, while accelerating hybrid and EV adoption
\end{itemize}
    
    \organization{Climate Damages Across the Globe \small{(Ongoing)}}{}{}
\begin{itemize}[leftmargin=2.5em, itemsep=-0.25em, topsep=0.5em]\vspace{-0.3em}
    \item Estimate heterogeneous country-level responses to temperature shocks using a global panel of economic and climatic data
    \item Combine machine learning methods to identify the economic and geoclimatic characteristics driving the heterogeneous effects
\end{itemize}

%%%%%%%%%%%%%%%%%%%%%%%%%%%%%%%%%%%%%%%%%%%%%%%%%%%%%%%%%%%%%%%%%%%%%%
% Work Experience
\sectionheading{Professional Experience}

    \organization{Pacific Investment Management Company (PIMCO)}{Newport Beach, United States}{Jun 2026 -- Aug 2026}
	\begin{position}{Quantitative Research Analyst}{}	
		\item Reviewed the literature studying AI and developed a HANK model to study its medium-run macroeconomic effects
        \item Developed an interpretable model to forecast aggregate S\&P 500 earnings per share, reducing the mean squared error by 18\% relative to a historical-mean benchmark and achieving 75\% directional accuracy   
	\end{position}

	\organization{London School of Economics}{London, United Kingdom}{Sep 2020 -- Sep 2021}
	\begin{position}{Research Assistant, }{Centre for Macroeconomics. Advisors: Profs. Benjamin Moll \& Ricardo Reis}	
		\item Solved and simulated continuous-time heterogeneous-agent models numerically
			%survey of professional forecasters, michigan survey of consumers
		\item Modeled rare inflation events using data-driven threshold selection and MLE of light- and heavy-tailed distributions
			%First define threshold for what an extreme-event is
			%Then fit fat-tailed distributions using either closed-form solutions or MLE
		\item Analyzed microdata from consumer and professional surveys to study the formation and distribution of inflation expectations
	\end{position}


	\organization{European Central Bank}{Frankfurt am Main, Germany}{Jun 2018 -- Sep 2020}
	\begin{position}{Analyst, }{Crisis Management - Quantitative Cluster}
			\item Contributed to a natural language early-warning system processing 1+ million daily observations from a global media database
		%google big query
		%use GDELT full text api to retrieve text that mentioned particular banks/countries 
		%Global Database of Events, Language and Tone
		%created sentiment score of sentences using dictionary of pre-selected words
		%get total article count by theme (banking for example)
		%frequency count of articles
			\item Analyzed liquidity provisions during stress events and the impact of banking policy on profitability
			\item Contributed to the impact assessment of Basel III on MREL requirements
	\end{position}
	\begin{position}{Trainee, }{Stress-Test Modeling}
		\item Applied regression discontinuity and difference-in-differences methods to supervisory banking panel data to estimate the effects of stress tests and capital requirements on bank risk-taking and profitability
		\item Contributed to the development of a credit risk assessment tool using historical loan portfolio data
	\end{position}
%\organization{Bank BPI}{Lisbon, Portugal}{Jun 2017 - Aug 2017}
%	\begin{position}[noitemize]{Intern, }{Private Banking}
%	\end{position}
	
%%%%%%%%%%%%%%%%%%%%%%%%%%%%%%%%%%%%%%%%%%%%%%%%%%%%%%%%%%%%%%%%%%%%%%
% Skills
\sectionheading{Tools \& Skills}
\vspace{-0.4em}

\noindent
\begin{minipage}[t]{0.48\textwidth}
\begin{itemize}[leftmargin=*, itemsep=0em, topsep=0.3em]
    \item \textbf{Programming:} Python, R, MATLAB
    \item \textbf{Time Series:} State-Space Models, Nonlinear Models, Stochastic Volatility, Bayesian Inference, Forecasting
    \item \textbf{Econometrics:} Shift-share IV, DiD, RDD, Statistical Learning
    \item \textbf{Computing:} High-performance computing (SLURM), Git, \LaTeX
\end{itemize}
\end{minipage}%
\hfill%
\begin{minipage}[t]{0.48\textwidth}
\begin{itemize}[leftmargin=0.25em, itemsep=0.25em, topsep=0.3em]
    \item \textbf{Quantitative Macro:} Heterogeneous-Agent Models, Continuous-Time Methods, Sequence-Space Jacobian, Dynare
    \item \textbf{Datasets:} Compustat, CRSP, CES, PSID, IBES, SPF, SCE, MER, SEDS, CMIP6
    \item \textbf{Languages:} English, Portuguese (fluent); French, Spanish (conversational)
\end{itemize}
\end{minipage}


		
\end{document}
